\documentclass[a4paper,12pt]{article}
\usepackage[T2A]{fontenc}
\usepackage[utf8]{inputenc}
\usepackage{cmap}
\usepackage[english, russian]{babel}
\usepackage{wrapfig}
\usepackage[backend=biber,style=gost-numeric]{biblatex} % стиль ГОСТ
\addbibresource{literatura.bib}
\usepackage{misccorr} % в заголовках появляется точка, но при ссылке на них ее нет
\usepackage{amssymb,amsfonts,amsmath,amsthm}  
\usepackage{indentfirst}
\usepackage[usenames,dvipsnames]{color} 
\usepackage[unicode,hidelinks]{hyperref}
% \hypersetup{%
%     pdfborder = {0 0 0}
\usepackage{multirow}
\usepackage{makecell,multirow} 
\usepackage{ulem}
%\usepackage{hyperref}
%\usepackage{graphicx,wrapfig}
%\renewcommand\epsilon{\varepsilon}
% \renewcommand\epsilon{\varepsilon}
%\graphicspath{{img/}}
\usepackage{geometry}
\geometry{left=1cm,right=1cm,top=2cm,bottom=2cm,bindingoffset=0cm,headheight=15pt}
\usepackage{fancyhdr} 
\linespread{1.2} 
\frenchspacing 
\renewcommand{\labelenumii}{\theenumii} 
% \usepackage{caption}
%%%%%%%%%%%%%%%%%%%%%%%%%%%%%%%%%%%%%%%%%%%%%%%%%%%%%%%%%%%%%%%%%%%%%%%%%%%%%%%
%%%%%%%%%%%%%%%%%%%%%%%%%%%%%%%%%%%%%%%%%%%%%%%%%%%%%%%%%%%%%%%%%%%%%%%%%%%%%%%

\def\labauthor{Сарафанов И.Г.}
\def\labauthors{\labauthor}
\def\labnumber{127}
\def\labtheme{Таблица производных и интегралов}

%%%%%%%%%%%%%%%%%%%%%%%%%%%%%%%%%%%%%%%%%%%%%%%%%%%%%%%%%%%%%%%%%%%%%%%%%%%%%%%
	%применим колонтитул к стилю страницы
\pagestyle{fancy} 
	%очистим "шапку" страницы
\fancyhead{} 
	%слева сверху на четных и справа на нечетных
\fancyhead[L]{\labauthors} 
	%справа сверху на четных и слева на нечетных
\fancyhead[R]{Погрешности для Лизы} 
	%очистим "подвал" страницы
\fancyfoot{} 
	% номер страницы в нижнем колинтуле в центре
\fancyfoot[C]{\thepage} 
%%%%%%%%%%%%%%%%%%%%%%%%%%%%%%%%%%%%%%%%%%%%%%%%%%%%%%%%%%%%%%%%%%%%%%%%%%%%%%%
\usepackage{float}
\usepackage[mode=buildnew]{standalone}
\usepackage{tikz} 
% \usepackage{subcaption}
\usepackage{tikz,csvsimple,physics}
\makeatother
%\usepackage{booktabs}
\usepackage{pgfplots, pgfplotstable}
\usepackage[outline]{contour}
\usepackage{tocloft}
\renewcommand{\cftsecleader}{\cftdotfill{\cftdotsep}}
%
%
%
%
%
\begin{document}

\section*{Обработка результатов измерений и расчёт погрешности углового ускорения \(\gamma_z\)}

\subsection*{1. Расчётная формула для углового ускорения}

В работе угловое ускорение маятника \(\gamma_z\) определяется косвенно — через измерение времени \(t\) опускания разгонного груза с высоты \(h\) при известном радиусе шкива \(r\). Используется формула, связывающая линейное ускорение груза с угловым ускорением маятника (кинематическая связь \(a = \gamma_z r\)):

\[
a = \frac{2h}{t^2}, \quad \gamma_z = \frac{a}{r} = \frac{2h}{r t^2}. \tag{1}
\]

Таким образом, \(\gamma_z\) является функцией трёх непосредственно измеряемых величин: высоты \(h\), времени \(t\) и радиуса шкива \(r\).

\subsection*{2. Расчёт погрешности косвенных измерений}

Поскольку \(\gamma_z\) не измеряется непосредственно, а вычисляется по формуле (1), её абсолютная погрешность находится методами теории погрешностей для косвенных измерений. В работе используются два подхода: квадратичный (статистический) и линейный (метод границ).

В общем виде для функции \(f = f(x_1, x_2, \dots, x_n)\) абсолютная погрешность рассчитывается с помощью частных производных, которые показывают чувствительность функции к изменению каждого аргумента.

\subsubsection*{2.1. Частные производные функции \(\gamma_z(h, r, t)\)}

Для применения любого из методов необходимо вычислить частные производные \(\gamma_z\) по каждому аргументу:

\begin{itemize}
    \item Производная по высоте \(h\):
    \[
    \frac{\partial \gamma_z}{\partial h} = \frac{\partial}{\partial h} \left( \frac{2h}{r t^2} \right) = \frac{2}{r t^2}.
    \]
    \item Производная по радиусу шкива \(r\):
    \[
    \frac{\partial \gamma_z}{\partial r} = \frac{\partial}{\partial r} \left( \frac{2h}{r t^2} \right) = -\frac{2h}{r^2 t^2}.
    \]
    \item Производная по времени \(t\):
    \[
    \frac{\partial \gamma_z}{\partial t} = \frac{\partial}{\partial t} \left( \frac{2h}{r} t^{-2} \right) = \frac{2h}{r} \cdot (-2) t^{-3} = -\frac{4h}{r t^3}.
    \]
\end{itemize}

Знак минус у производных по \(r\) и \(t\) указывает на то, что увеличение этих величин приводит к уменьшению \(\gamma_z\). При расчёте погрешности методом наихудшего случая знак не важен — берётся модуль.
\[
\Delta \gamma_z = \left| \frac{2}{r t^2} \right| \Delta h + 
\left| -\frac{2h}{r^2 t^2} \right| \Delta r + 
\left| -\frac{4h}{r t^3} \right| \Delta t.
\]

Так как модуль отрицательного числа даёт положительное, окончательно:

\[
\boxed{\Delta \gamma_z = \frac{2}{r t^2} \Delta h + \frac{2h}{r^2 t^2} \Delta r + \frac{4h}{r t^3} \Delta t}. \tag{3}
\]

\textbf{Значения \textit{абсолютных}(маркер - $\Delta$) погрешностей $\Delta t, \Delta h, \Delta r$ нам давались в протоколе} :))
\subsubsection*{2.2. Квадратичный (статистический) метод, (оно наврное не ненужно)}

Этот метод является более корректным, так как учитывает вероятностный характер случайных погрешностей и предполагает их независимость. Абсолютная погрешность \(\Delta \gamma_z\) рассчитывается как корень квадратный из суммы квадратов вкладов каждой частной погрешности:

\[
\Delta \gamma_z = \sqrt{ \left( \frac{\partial \gamma_z}{\partial h} \Delta h \right)^2 + \left( \frac{\partial \gamma_z}{\partial r} \Delta r \right)^2 + \left( \frac{\partial \gamma_z}{\partial t} \Delta t \right)^2 }.
\]

Подставляя выражения для частных производных, получаем рабочую формулу:

\[
\boxed{\Delta \gamma_z = \sqrt{ \left( \frac{2}{r t^2} \Delta h \right)^2 + \left( -\frac{2h}{r^2 t^2} \Delta r \right)^2 + \left( -\frac{4h}{r t^3} \Delta t \right)^2 }}.
\]

Так как квадрат отрицательного числа положителен, знаки минус можно опустить:

\[
\Delta \gamma_z = \sqrt{ \left( \frac{2 \Delta h}{r t^2} \right)^2 + \left( \frac{2h \Delta r}{r^2 t^2} \right)^2 + \left( \frac{4h \Delta t}{r t^3} \right)^2 }. \tag{2}
\]

\subsubsection*{2.3. Линейный метод (метод наихудшего случая или метод границ)}

Этот метод даёт оценку погрешности «сверху», предполагая, что все частные погрешности складываются в одном направлении (наихудшее сочетание знаков). Абсолютная погрешность находится как сумма модулей вкладов:

\[
\Delta \gamma_z = \left| \frac{\partial \gamma_z}{\partial h} \right| \Delta h + \left| \frac{\partial \gamma_z}{\partial r} \right| \Delta r + \left| \frac{\partial \gamma_z}{\partial t} \right| \Delta t.
\]

Подставляя модули частных производных, получаем:

\[
\boxed{\Delta \gamma_z = \frac{2}{r t^2} \Delta h + \frac{2h}{r^2 t^2} \Delta r + \frac{4h}{r t^3} \Delta t}. \tag{3}
\]

\subsection*{3. Рекомендации по применению}

\begin{itemize}
    \item Для окончательной оценки погрешности в лабораторных работах чаще используется **квадратичный метод** (формула (2)), так как он более реалистично отражает влияние случайных ошибок.
    \item **Линейный метод** (формула (3)) удобен для быстрых оценок или когда погрешности имеют систематический характер и их знаки известны. Он всегда даёт большую (более грубую) оценку по сравнению с квадратичным.
    \item Погрешности прямых измерений \(\Delta h\), \(\Delta r\) и \(\Delta t\) определяются по классу точности приборов или как приборная погрешность (для \(h\) и \(r\)), либо как случайная погрешность времени (по разбросу измерений \(t\)).
\end{itemize}
Абсолютная погрешность для массы $\Delta m = n \cdot \Delta m_1 $
\end{document}